% !TeX root = ../main.tex

\chapter{案例分析}
\label{cha:case_analysis}

本章从自建集评测结果中选取5个成功案例和5个失败案例进行深度分析，旨在揭示模型的能力边界和典型失败模式，为系统改进提供方向性指导。

\section{成功案例分析}

\subsection{案例一：猫的物体识别}

\begin{itemize}
  \item \textbf{图片类型}：自然场景照片
  \item \textbf{问题}：请描述这张图片中的主要内容。
  \item \textbf{期望答案}：图片中有一只猫。
  \item \textbf{模型回答}：一张小猫的特写照片，主体是一只年幼的虎斑猫，它正趴在一块浅绿色的绒面织物上。小猫拥有明亮蓝绿色眼睛，瞳孔清晰，鼻子是粉嫩的肉色，嘴周和下巴有白色毛发，典型的虎斑猫纹路——头部、耳朵和背部覆盖着深灰与黑色相间的条纹。
  \item \textbf{分析}：模型不仅准确识别了"猫"这一核心对象，还提供了丰富的细节描述（品种"虎斑猫"、颜色"蓝绿色眼睛""粉嫩肉色鼻子"、姿态"趴"、环境"浅绿色绒面织物"）。这表明模型在自然场景下的物体识别和属性提取能力较为成熟，生成的回答内容远超期望答案的简洁程度，属于高质量的成功案例。
\end{itemize}

\subsection{案例二：秋日风景描述}

\begin{itemize}
  \item \textbf{图片类型}：自然场景照片
  \item \textbf{问题}：这张图片展现了什么样的风景？
  \item \textbf{期望答案}：图片中有一条道路和秋天的树木。
  \item \textbf{模型回答}：一幅绚烂迷人的秋日森林小径风景。整幅画面被秋天的主色调——炽烈的红色、橙色、金黄与赭石色所笼罩。枫树等落叶阔叶树的叶子红得耀眼，地面铺满了层层叠叠的落叶，与蜿蜒的小路交织成一幅暖色调的地毯。空间层次丰富，一条浅色的土路或碎石小径从前景向远处蜿蜒伸展。
  \item \textbf{分析}：模型准确抓住了"秋天的树木"和"道路"两个核心要素，并以高质量的中文描写扩展了场景的氛围和细节。该案例体现了模型在自然场景理解中对色彩、空间和季节特征的综合把握能力，回答的文学性和信息量均达到较高水平。
\end{itemize}

\subsection{案例三：白色雏菊识别}

\begin{itemize}
  \item \textbf{图片类型}：自然场景照片（花卉特写）
  \item \textbf{问题}：图中有什么植物？
  \item \textbf{期望答案}：图中有一朵白色的花。
  \item \textbf{模型回答}：图中植物是雏菊（\textit{Leucanthemum vulgare}）。特征包括：白色舌状花瓣呈放射状排列，中央为黄色管状花组成的花盘，花茎细长、绿色、略带弯曲，整体形态简洁清新，是典型的菊科植物。
  \item \textbf{分析}：模型超越了"白色花"这一基本要求，提供了科属级别的精确识别（菊科雏菊属）甚至拉丁学名。该案例说明模型在常见花卉的细粒度识别上具有超出预期的知识储备，能够将视觉特征与植物学知识进行有效关联。
\end{itemize}

\subsection{案例四：文档OCR——文字识别}

\begin{itemize}
  \item \textbf{图片类型}：文档截图
  \item \textbf{问题}：文档的标题是什么？
  \item \textbf{期望答案}：实验报告。
  \item \textbf{模型回答}：文档的标题是"实验报告"。
  \item \textbf{分析}：模型简洁准确地完成了OCR识别任务。文档类问题的答案通常直接可见于图片中，OCR是Qwen3-VL-Flash的强项，响应速度快（0.77s）。该案例代表了文档OCR子类别中63.0\%准确率的典型成功场景——当文字清晰、排版规范时，模型表现稳定可靠。
\end{itemize}

\subsection{案例五：场景推理——室内外判断}

\begin{itemize}
  \item \textbf{图片类型}：自然场景照片（户外湖泊）
  \item \textbf{问题}：这个场景是在室内还是室外？请给出理由。
  \item \textbf{期望答案}：这个场景是在室外。
  \item \textbf{模型回答}：这个场景是在室外。理由如下：画面中是一片开阔平静的湖面，水面倒映着小船和人，说明是自然水体；天空灰蒙蒙、有雾，是典型的户外天气现象；两人坐在独木舟中，穿着防风保暖的连帽外套，符合户外水上活动的着装。
  \item \textbf{分析}：在推理判断类问题上，模型不仅给出了正确结论，还提供了"开阔水面""雾气天气""户外着装"三条具体推理依据，展示了从视觉特征到逻辑推理的链条式思维过程。
\end{itemize}

\section{失败案例分析}

\subsection{案例一：否定回答与期望矛盾——交通工具识别}

\begin{itemize}
  \item \textbf{图片类型}：自然场景（城市天际线远景）
  \item \textbf{问题}：图中有什么交通工具？
  \item \textbf{期望答案}：图中有一辆自行车。
  \item \textbf{模型回答}：图中没有可见的交通工具。这张照片展示的是纽约市天际线的远景，画面主体是密集的城市建筑群。在图像的前景和中景，能看到大量建筑物的屋顶、烟囱、天线等，但没有清晰可辨的汽车、公交车、地铁、飞机、船只或其他交通工具。
  \item \textbf{失败原因}：评测集中期望答案"图中有一辆自行车"与实际图片内容不匹配（该图确实为城市天际线远景，无可辨识的自行车）。模型做出了正确的视觉判断（"没有可见的交通工具"），但在评测框架中被判为错误。这属于\textbf{评测集设计问题}——部分问题预设了不存在的物体。
\end{itemize}

\subsection{案例二：否定回答与期望矛盾——建筑识别}

\begin{itemize}
  \item \textbf{图片类型}：自然场景（海岸礁石）
  \item \textbf{问题}：图中有什么建筑？
  \item \textbf{期望答案}：图中有城市建筑。
  \item \textbf{模型回答}：图中没有建筑。画面展示的是自然海岸景观：巨大的海浪猛烈撞击着深色湿润的岩石，激起大量白色浪花和水雾。背景是淡蓝色的天空和远处的海面，整体呈现的是海洋与岩石的动态互动，属于自然风光，未见任何人工建筑。
  \item \textbf{失败原因}：与案例一类似，图片本身是纯粹的自然海岸景观，不存在人工建筑。模型对场景的判断完全正确，但评测集中期望答案为肯定型。"自然海岸景观"与"城市建筑"的根本矛盾导致评测失败。该案例揭示了评测算法的一个系统性局限——无法区分"模型判断错误"和"标注答案与图片事实不符"。
\end{itemize}

\subsection{案例三：否定回答与期望矛盾——动物识别}

\begin{itemize}
  \item \textbf{图片类型}：自然场景（人物剪影+日落）
  \item \textbf{问题}：图中有什么动物？
  \item \textbf{期望答案}：图中有一只狗。
  \item \textbf{模型回答}：图中没有动物。画面主体是一位金发女性背对镜头坐在岩石或堤岸上，面向远方的湖面和树林。背景是日落或日出时分的自然风光。整个场景中只有人物、自然环境，未见任何可辨识的动物。
  \item \textbf{失败原因}：图片焦点是人物剪影和日落风光，不存在动物。模型如实描述了所见内容。评测集中"图中有一只狗"的标注与图片实际内容存在差异。
\end{itemize}

\subsection{案例四：否定回答与期望矛盾——花卉识别}

\begin{itemize}
  \item \textbf{图片类型}：自然场景（织物特写）
  \item \textbf{问题}：图中有什么花卉？
  \item \textbf{期望答案}：图中有向日葵。
  \item \textbf{模型回答}：图中没有花卉。这是一张织物的特写照片，主要展示了一条色彩丰富的针织围巾/披肩，图案为几何抽象风格，包含锐利的三角形、锯齿形等图形，色彩以粉红、紫红、青绿、深蓝、赭石等为主。没有任何可辨识的花卉图案。
  \item \textbf{失败原因}：模型正确地识别了图片为纺织品而非花卉照片。该案例再次说明否定回答场景下的评测困境——当问题假设图片中存在某类物体而实际不存在时，模型的最佳回答反而是"没有"。
\end{itemize}

\subsection{案例五：时钟读数识别精度不足}

\begin{itemize}
  \item \textbf{图片类型}：自然场景（多时钟展示）
  \item \textbf{问题}：顶部时钟显示的是几点？（TextVQA公开集案例）
  \item \textbf{期望答案}：10:07
  \item \textbf{模型回答}：顶部时钟显示：时针指向刚过10的位置，分针指向1的位置，因此大约是10:05。
  \item \textbf{失败原因}：这是一个细粒度视觉识别失败的典型案例。模型正确描述了指针位置，但读取精度不足：将分针在7分钟位置误判为5分钟位置。模拟时钟的精确读数需要毫米级的指针角度识别能力，当前VLM在此类高精度视觉任务上仍有改进空间。这也反映了TextVQA评测中字符串严格匹配的局限性——模型回答"10:05"与期望"10:07"仅差2分钟，语义上已非常接近。
\end{itemize}

\section{失败模式总结}

综合5个失败案例的分析，可以归纳出vlm-chat在当前阶段的\textbf{三类主要失败模式}：

\begin{enumerate}
  \item \textbf{否定回答冲突（案例一至四）}：模型对"图中不存在某物"的正确判断与评测集的肯定型期望答案矛盾。这类失败在所有失败案例中占比最高（自建集中约占60\%），反映了评测设计的系统性问题而非模型能力的不足；

  \item \textbf{细粒度视觉识别不足（案例五）}：在需要精确读取刻度、角度等信息的场景下，模型的视觉精度有限。模拟时钟读数、仪表盘识别等任务需要亚像素级的定位能力；

  \item \textbf{评测算法局限性}：当前基于jieba分词的匹配算法对否定句、同义表达、部分正确等场景的判别能力有限。例如案例五中"10:05"与"10:07"的2分钟误差在语义层面可接受，但在字符串匹配层面被严格判错。
\end{enumerate}

这些失败模式的识别为下一章的反思与改进讨论提供了实证基础。
